% Cover Letter — CRANE-Z — target: Nature Aging
% Prepared 2026-08-12 for top-journal submission package.

\textbf{To:} The Editor-in-Chief, \textit{Nature Aging}

\medskip

\noindent\textbf{Re:} Submission of \textit{CRANE-Z: dual-modality deep learning reveals decelerated biological aging in long-lived individuals across three independent cohorts, two ethnicities and two measurement platforms}

\medskip

\noindent\textbf{Manuscript type:} Article (Original Research)

\medskip

\noindent\textbf{Corresponding author:} Xuepeng Fu, Qiqihar University, \href{mailto:02383@qqhru.edu.cn}{02383@qqhru.edu.cn}

\medskip

\noindent\textbf{Lead contact:} Tongfei Shen, Qingdao Huanghai University, \href{mailto:02383@qqhru.edu.cn}{02383@qqhru.edu.cn}

\medskip

\hrule

\medskip

Dear Editors,

We report CRANE-Z, a dual-modality deep classifier that integrates a gene-module Transformer with an LM22 immune deconvolution branch and mask-based feature pretraining, and demonstrate that it learns a transferable longevity signature across three fully independent cohorts, two ethnicities and two measurement technologies. The work, we believe, makes three advances directly relevant to \textit{Nature Aging}'s scope.

\textbf{First, a new discriminator that is not a clock.} Existing longevity transcriptomic classifiers are descriptive, they estimate biological age but do not frame the problem as cross-cohort phenotype discrimination. CRANE-Z is trained as a 5-fold by 5-seed discriminator of long-lived (RLL, $\geq$70 y) versus young-control (RYC, $\leq$50 y) phenotypes, and is evaluated externally on tissues and ethnicities that the model never saw at training. On the internal Chinese longevity cohort (1,715 blood transcriptomes), long-lived individuals carry a mean aging-acceleration of $-8.1$ years (group separation 14.69 years, $p<10^{-22}$, Cohen $d=1.18$), and the model age ranks individuals with a Spearman correlation of 0.816 with chronological age.

\textbf{Second, three-cohort, two-ethnicity, two-platform replication of the decelerated biological-clock signal.} CRANE-Z generalizes to GTEx whole blood (n=64) and skeletal muscle (n=58) with the same negative-aging-acceleration direction in long-lived individuals ($p=0.001$), and replicates on a fully independent Israeli whole-blood cohort profiled on the Affymetrix PrimeView microarray (GSE123696, n=66), where the long-lived group shows a mean aging-acceleration of $-3.11$ years versus $+14.78$ years in young controls ($\Delta=-17.89$ years, Welch $p=5.9\times10^{-6}$, Mann-Whitney $p=2.9\times10^{-5}$, Cohen $d=1.39$, 5/5 fold direction-consistent, Spearman $\rho=-0.569$, $p=4\times10^{-7}$). This three-cohort evidence chain is, to our knowledge, the largest cross-platform validation of any transcriptomic longevity classifier.

\textbf{Third, a defined molecular mechanism for the signature.} Pre-ranked gene set enrichment analysis on the genome-wide gene-age rank list (16,197 genes, 1,000 permutations, four pathway libraries: KEGG, GO Biological Process, Reactome, WikiPathways) shows that the age axis CRANE-Z learns is dominated by immune-effector biology. The strongest hits are antigen processing and presentation (KEGG, NES=2.30, FDR$<10^{-3}$), regulation of dendritic cell differentiation (NES=2.62, FDR=0.013), natural-killer-cell-mediated cytotoxicity (NES=2.39, FDR=0.031), NK chemotaxis regulation (NES=2.36, FDR=0.028), cellular response to interferon-$\gamma$ (NES=2.36, FDR=0.032) and negative regulation of interleukin-1$\beta$ production (NES=2.40, FDR=0.034), with 18 of 19 FDR$<0.05$ hits age-up enriched. This places the aging-acceleration readout on a defined molecular substrate and mechanistically explains the cross-cohort immune signature of resting NK elevation and naive B reduction that the model recovers. The pathway convergence across four libraries makes the mechanism claim robust to library choice.

The cohort sizes are consistent with the field and the long-lived group definition follows the well-established RYC/RLL phenotype construction. The training code, the trained model weights, the processed tensors used in this work, the per-sample predictions underlying every reported table and figure, and the regeneration scripts are all freely available at \href{https://github.com/SHENTongfei/CRANE-Z}{https://github.com/SHENTongfei/CRANE-Z} (MIT license). The data sources are public. We are not aware of any work that has achieved the same combination of cross-cohort, cross-tissue, cross-ethnicity and cross-platform replication of a decelerated biological-age signal with a defined molecular-mechanism readout. We therefore believe the work is a strong fit for \textit{Nature Aging} and would be honoured to have it considered for peer review.

The manuscript has not been published and is not under consideration elsewhere. All authors have approved the submission and declare no competing interests.

Yours sincerely,

\medskip

\noindent Tongfei Shen (lead contact), on behalf of all co-authors

\medskip

\noindent Qiqihar University / Qingdao Huanghai University, 2026-08-12
